\chapter{Penser en espérance mathématique}\label{ch:m2:thinking-in-expected-value}

En 2016, deux chercheurs ont remis à 61 jeunes professionnels de la finance et étudiants 25 dollars
et 30 minutes pour parier sur une pièce qui, leur avait-on dit, tombait sur pile dans 60~\% des cas. Ils
pouvaient miser n'importe quel montant sur l'un ou l'autre côté, aussi souvent qu'ils le voulaient, dans la
limite d'un gain maximal non divulgué de 250 dollars. Une mise fixe de 10 à 20~\% du capital sur pile
aurait permis d'atteindre ce maximum dans 95~\% des cas. Seuls 21~\% des joueurs y sont
parvenus~; un tiers a fini avec moins qu'il n'avait au départ, et 28~\% ont fait faillite, perdant
tout sur une pièce biaisée en leur faveur. Ils ont trop misé, joué de façon erratique, et certains ont misé sur face.
Ce chapitre traite de l'arithmétique qu'ils ont négligée~: l'edge et la \omterm{def:m2:thinking-in-expected-value:edge}{cote}, la fraction de la
richesse qui maximise la croissance à long terme, les raisons pour lesquelles les traders misent moins
que cela, et de combien moins lorsque l'edge lui-même n'est qu'une estimation.

\section{Edge et cote}

\begin{definition}[Edge, cote]\label{def:m2:thinking-in-expected-value:edge}
La \emph{cote}\index{cote} $b$ d'un pari est ce qu'il rapporte par unité misée s'il est gagnant,
la mise étant perdue dans le cas contraire. Son \emph{edge}\index{edge} est le profit espéré par
unité misée, soit $pb - q$, où $p$ est la probabilité de gagner et $q = 1 - p$.
\end{definition}

La pièce paie la mise à l'égalité, $b = 1$, et le pari gagne avec $p = 0.6$~: son edge vaut $0.6 - 0.4 = 0.2$,
soit 20 centimes par dollar misé. Un edge positif indique que le pari vaut la peine d'être pris~; il ne dit rien
sur la mise à engager. Si l'on mise tout à chaque fois, un seul face met fin au jeu~:
l'espérance de chaque pari est positive, mais la probabilité de survivre à 300 lancers
vaut $0.6^{300}$. La question n'est pas l'edge d'un pari isolé, mais l'effet, sur la richesse, de miser
une fraction donnée de celle-ci, reconduite pari après pari.

\begin{definition}[Taux de croissance]\label{def:m2:thinking-in-expected-value:growth}
Le \emph{taux de croissance}\index{taux de croissance} d'une règle de mise est l'espérance du logarithme
du rapport entre la richesse après un pari et la richesse avant celui-ci. En misant une fraction $f$ de la richesse sur
un pari de \omterm{def:m2:thinking-in-expected-value:edge}{cote} $b$ et de probabilité de gain $p$, il vaut
\[
g(f) = p\ln(1 + fb) + q\ln(1 - f).
\]
\end{definition}

Après $n$ paris, le logarithme de la richesse est une somme de $n$ termes indépendants~; par conséquent, selon la loi
des grands nombres, la richesse croît presque sûrement comme $e^{n g(f)}$~: c'est le \omterm{def:m2:thinking-in-expected-value:growth}{taux de croissance}, et non l'edge,
qui détermine où aboutit la richesse. L'espérance de la richesse après $n$ paris vaut $(1 + f(pb -
q))^n$, qui croît sans limite avec $f$~; la médiane, elle, au plus $e^{n g(f)}$, ne suit pas cette croissance. Pour
$f = 0.2$ sur la pièce, le gain espéré est de 4~\% par lancer, et l'espérance de la valeur de 300 lancers
partant de 25 dollars vaut 3.2~millions de dollars, chiffre cité par Haghani et Dewey~; la médiane est
10\,504~dollars.

\section{Le critère de Kelly}

\begin{definition}[Critère de Kelly]\label{def:m2:thinking-in-expected-value:kelly}
Le \emph{critère de Kelly}\index{critère de Kelly} consiste à miser, à chaque pari, la fraction de
la richesse qui maximise le \omterm{def:m2:thinking-in-expected-value:growth}{taux de croissance}.
\end{definition}

\begin{proposition}[La fraction de Kelly]\label{prop:m2:thinking-in-expected-value:kelly}
Pour un pari de \omterm{def:m2:thinking-in-expected-value:edge}{cote} $b$ et de probabilité de gain $p$, le \omterm{def:m2:thinking-in-expected-value:growth}{taux de croissance} est maximal en
\[
f^* = p - \frac{q}{b} = \frac{\text{edge}}{\text{cote}},
\]
si l'edge est positif, et en $f^* = 0$ sinon. Pour de petits paris de moyenne $\mu$ et
de variance $s^2$ par unité misée, $g(f) \approx f\mu - \tfrac12 f^2 s^2$, maximal en $f^* =
\mu / s^2$, avec $g(f^*) = \mu^2 / (2s^2)$.
\end{proposition}

\begin{proof}
$g'(f) = pb/(1 + fb) - q/(1 - f)$ s'annule en $f = (pb - q)/b$, et $g$ est concave. Pour la
seconde forme, on développe le logarithme à l'ordre deux~: $\ln(1 + fX) \approx fX - \tfrac12 f^2X^2$
et on prend l'espérance.
\end{proof}

La fraction de Kelly de la pièce est 0.2~: il faut miser 20~\% du capital sur pile à chaque fois, pour un
\omterm{def:m2:thinking-in-expected-value:growth}{taux de croissance} de 2.01~\% par lancer. La forme quadratique montre la nature du compromis
(\cref{fig:m2:thinking-in-expected-value:growth})~: la croissance s'élève, atteint son maximum en $f^*$, puis retombe
à zéro vers le double de $f^*$~; au-delà, plus on mise, plus la ruine est rapide,
quelle que soit la taille de l'edge. Pour la pièce, le \omterm{def:m2:thinking-in-expected-value:growth}{taux de croissance} s'annule pour une mise de 38.9~\%.

\begin{omfigure}
\begin{tikzpicture}
\begin{axis}[width=11cm,height=4.8cm,
  xlabel={fraction de la richesse misée à chaque lancer, $f$},ylabel={taux de croissance (\% par lancer)},
  tick label style={font=\small},label style={font=\small},
  xmin=0,xmax=0.46,ymin=-1,ymax=2.8,grid=major,grid style={gray!25},
  xticklabel style={/pgf/number format/.cd,fixed,precision=2},
  table/col sep=comma]
\addplot[omDef,very thick] table[x=f,y=g]{figdata/markets-2/29-thinking-in-expected-value/growth.csv};
\addplot[gray,thin] coordinates {(0,0) (0.46,0)};
\addplot[only marks,mark=*,mark size=1.8pt,black] coordinates {(0.2,2.0136) (0.3894,0)};
\node[font=\footnotesize,anchor=south] at (axis cs:0.2,2.05) {Kelly~: 0.2};
\node[font=\footnotesize,anchor=east] at (axis cs:0.375,-0.45) {croissance nulle~: 0.389};
\end{axis}
\end{tikzpicture}

\omcaption{\omterm{def:m2:thinking-in-expected-value:growth}{Taux de croissance} de la richesse par lancer de la pièce à 60~\% à égalité de \omterm{def:m2:thinking-in-expected-value:edge}{cote}, en fonction de la
fraction de richesse misée. Il atteint son maximum à la fraction de Kelly, 0.2, avec 2.01~\% par lancer, et
s'annule à 0.389, soit environ le double de Kelly. Données~: le tutoriel du chapitre.}\label{fig:m2:thinking-in-expected-value:growth}
\end{omfigure}

\section{Kelly fractionnaire et risque de ruine}

\begin{definition}[Kelly fractionnaire, drawdown, risque de ruine]\label{def:m2:thinking-in-expected-value:frac}
Le \emph{Kelly fractionnaire}\index{Kelly fractionnaire} consiste à miser un multiple $c < 1$ de la
fraction de Kelly. Un \emph{drawdown}\index{drawdown} est une baisse de la richesse par rapport à son précédent sommet,
mesurée en fraction de ce sommet. Le \emph{risque de ruine}\index{risque de ruine} d'une règle de mise est
la probabilité que la richesse tombe, à un moment donné, à une fraction donnée de sa
valeur initiale.
\end{definition}

Miser selon le \omterm{def:m2:thinking-in-expected-value:kelly}{critère de Kelly} maximise la croissance, mais c'est une trajectoire violente~: la distribution de la richesse
après 300 lancers est extrêmement dispersée (\cref{fig:m2:thinking-in-expected-value:coin}).
Le demi-Kelly renonce à un quart du \omterm{def:m2:thinking-in-expected-value:growth}{taux de croissance} dans l'approximation quadratique, $g(cf^*)
= (2c - c^2)g(f^*)$, et une grande partie de la dispersion~: pour la pièce, la médiane après 300 lancers est
de 2\,279~dollars contre 10\,504, mais le dixième des pires issues est meilleur, 251 contre 121.
Miser le double de Kelly donne, à l'ordre deux, une croissance nulle~: la moitié des joueurs finissent sous leur mise de départ.

\begin{omfigure}
\begin{tikzpicture}
\begin{axis}[width=11cm,height=5cm,
  ylabel={richesse après 300 lancers (USD)},ymode=log,
  tick label style={font=\small},label style={font=\small},
  xtick={0,1,2,3},xticklabels={$\tfrac14$ de Kelly,$\tfrac12$ de Kelly,Kelly,double Kelly},
  xmin=-0.5,xmax=3.5,ymin=1e-4,ymax=5e6,grid=major,grid style={gray!25},
  legend columns=3,legend style={font=\footnotesize,at={(0.5,-0.25)},anchor=north,draw=none,fill=none,/tikz/every even column/.append style={column sep=8pt}},
  table/col sep=comma]
\addplot[omProp,only marks,mark=triangle*,mark size=2.6pt] table[x=k,y=p90]{figdata/markets-2/29-thinking-in-expected-value/coin.csv};
\addplot[omDef,only marks,mark=*,mark size=2.4pt] table[x=k,y=median]{figdata/markets-2/29-thinking-in-expected-value/coin.csv};
\addplot[omThm,only marks,mark=triangle*,mark options={rotate=180},mark size=2.6pt] table[x=k,y=p10]{figdata/markets-2/29-thinking-in-expected-value/coin.csv};
\addplot[gray,dashed,forget plot] coordinates {(-0.5,25) (3.5,25)};
\legend{90\up{e}~centile,médiane,10\up{e}~centile}
\end{axis}
\end{tikzpicture}

\omcaption{Richesse après 300 lancers de la pièce à 60~\% partant de 25~USD (en tirets), sans plafond, pour quatre
multiples de la fraction de Kelly, sur échelle logarithmique~: les centiles 10\up{e}, 50\up{e} et 90\up{e} de
4\,000 joueurs simulés pour chacun. Kelly a la médiane la plus élevée~; le demi-Kelly, le meilleur
dixième des pires issues~; le double de Kelly donne une médiane inférieure à la mise. Données~: le tutoriel du chapitre,
avec graine fixée.}\label{fig:m2:thinking-in-expected-value:coin}
\end{omfigure}

\begin{proposition}[Risque de chute d'une fraction]\label{prop:m2:thinking-in-expected-value:ruin}
Dans la limite en temps continu de petits paris, misant $c$ fois la fraction de Kelly, la
probabilité que la richesse tombe un jour à une fraction $\alpha$ de sa valeur initiale est
\[
P = \alpha^{2/c - 1}.
\]
À Kelly plein, la probabilité de tomber un jour à la moitié de sa richesse initiale vaut un demi.
\end{proposition}

\begin{proof}
Avec une mise $c\mu/s^2$, le logarithme de la richesse est un mouvement brownien de dérive $m = (c - \tfrac12
c^2)\mu^2/s^2$ et de variance $v = c^2\mu^2/s^2$ par unité de temps. Un mouvement brownien de dérive
$m > 0$ et de variance $v$ tombe un jour d'une quantité $d$ avec probabilité $e^{-2md/v}$~; avec $d = -\ln
\alpha$, $2m/v = 2/c - 1$.
\end{proof}

La formule se vérifie en simulation (la \cref{fig:m2:thinking-in-expected-value:dd})~: sur un pari
à 55~\% à mise égale sur 4\,000 paris, la part des trajectoires qui ont un jour été réduites de moitié vaut 48.4~\% à Kelly plein
contre 50~\% en théorie, 11.1~\% au demi-Kelly contre 12.5~\%. Inverser la formule donne le plus grand multiple de Kelly
pour une tolérance donnée~: pour que la probabilité de tomber un jour à la moitié reste
sous 10~\%, $c = 2/(1 + \ln 0.1/\ln 0.5) = 0.46$. Une chute depuis le sommet, que mesure le \omterm{def:m2:thinking-in-expected-value:frac}{drawdown}, se reproduit
sur un horizon assez long, quelle que soit la fraction~: le \omterm{def:m2:thinking-in-expected-value:frac}{risque de ruine} est
la version qui admet une réponse.

\begin{omfigure}
\begin{tikzpicture}
\begin{axis}[width=11cm,height=4.8cm,
  xlabel={multiple de la fraction de Kelly, $c$},ylabel={P(richesse divisée par 2)},
  tick label style={font=\small},label style={font=\small},
  xmin=0.2,xmax=1.55,ymin=0,ymax=0.85,grid=major,grid style={gray!25},
  yticklabel style={/pgf/number format/.cd,fixed,precision=1},
  legend columns=2,legend style={font=\footnotesize,at={(0.5,-0.3)},anchor=north,draw=none,fill=none,/tikz/every even column/.append style={column sep=8pt}},
  table/col sep=comma]
\addplot[omDef,very thick,domain=0.2:1.55,samples=60] {0.5^(2/x - 1)};
\addplot[omThm,only marks,mark=*,mark size=2pt] table[x=c,y=sim]{figdata/markets-2/29-thinking-in-expected-value/drawdown.csv};
\legend{$\alpha^{2/c-1}$ avec $\alpha = \tfrac12$,simulé (1\,000 trajectoires)}
\end{axis}
\end{tikzpicture}

\omcaption{Probabilité que la richesse tombe un jour à la moitié de sa valeur initiale, en fonction du
multiple de Kelly misé, sur un pari à 55~\% à mise égale~: la formule en temps continu et la
part de 1\,000 trajectoires simulées de 4\,000 paris. Données~: le tutoriel du chapitre,
avec graine fixée.}\label{fig:m2:thinking-in-expected-value:dd}
\end{omfigure}

\section{Quand l'edge lui-même est incertain}

Tout ce qui précède supposait l'edge connu. Un trader l'estime à partir d'un backtest ou de
l'historique propre de la stratégie, avec une erreur~; et la fraction de Kelly étant proportionnelle à
l'edge, si bien qu'un edge surestimé conduit à sur-miser, ce qui coûte plus de croissance
qu'une sous-mise de même ampleur.

\begin{proposition}[Réduction pour l'erreur d'estimation]\label{prop:m2:thinking-in-expected-value:shrink}
Supposons que la moyenne $\mu$ par unité misée soit estimée, à partir de $n$ transactions, par $\hat\mu$, avec
erreur type $s/\sqrt n$, et que le trader mise $c\hat\mu/s^2$. Dans l'approximation
quadratique, le \omterm{def:m2:thinking-in-expected-value:growth}{taux de croissance} espéré est maximal pour
\[
c^* = \frac{\mu^2}{\mu^2 + s^2/n} = \frac{n\,\mathrm{SR}^2}{n\,\mathrm{SR}^2 + 1},
\]
où $\mathrm{SR} = \mu/s$ est le ratio de Sharpe par transaction~; $n\,\mathrm{SR}^2$ est le carré
de la $t$-statistique de l'edge.
\end{proposition}

\begin{proof}
$E[g] = E[f]\mu - \tfrac12 E[f^2]s^2$ avec $f = c\hat\mu/s^2$, $E[\hat\mu] = \mu$ et
$E[\hat\mu^2] = \mu^2 + s^2/n$, donc $E[g] = c\mu^2/s^2 - \tfrac12 c^2(\mu^2 + s^2/n)/s^2$,
maximisée pour le $c$ énoncé.
\end{proof}

Une stratégie dont l'edge présente une $t$-statistique de 1, à peine distinguable de zéro,
doit miser la moitié de sa fraction de Kelly estimée~; une autre, avec une $t$-statistique de 3, neuf
dixièmes. Miser une fraction de Kelly relève, en partie, de cette arithmétique
(la \cref{fig:m2:thinking-in-expected-value:uncertain}).

\begin{omfigure}
\begin{tikzpicture}
\begin{axis}[width=11cm,height=4.8cm,
  xlabel={multiple de la fraction de Kelly estimée, $c$},ylabel={croissance (pb par transaction)},
  tick label style={font=\small},label style={font=\small},
  xmin=0,xmax=1.6,ymin=-20,ymax=45,grid=major,grid style={gray!25},
  table/col sep=comma]
\addplot[omDef,very thick,mark=*,mark size=1.4pt] table[x=c,y=g]{figdata/markets-2/29-thinking-in-expected-value/uncertain.csv};
\addplot[omThm,thick,dashed] coordinates {(0.669,-20) (0.669,45)};
\addplot[gray,thin] coordinates {(0,0) (1.6,0)};
\node[font=\footnotesize,anchor=north west] at (axis cs:0.69,-5) {$c^* = 0.669$};
\end{axis}
\end{tikzpicture}

\omcaption{Croissance par transaction, en points de base, d'un trader qui estime le taux de gain d'un
pari à \omterm{def:m2:thinking-in-expected-value:edge}{cote} 1 contre 1 gagnant dans 55~\% des cas, à partir de 200 transactions, puis mise $c$ fois la fraction de Kelly
estimée pendant 500 transactions, moyennée sur 3\,000 répétitions. Le meilleur multiple de la
simulation est 0.7~; la formule donne 0.669. Miser la fraction de Kelly estimée en totalité
rapporte 19~\% de moins. Données~: le tutoriel du chapitre, avec graine fixée.}\label{fig:m2:thinking-in-expected-value:uncertain}
\end{omfigure}

L'approximation quadratique masque aussi les queues épaisses~: une transaction qui peut perdre plus que
sa mise, ou une distribution des rendements dont les valeurs extrêmes sont sous-estimées, appelle
des fractions encore plus petites. Et Kelly maximise la croissance d'un seul capital, sans
contrainte extérieure~; un fonds dont les investisseurs se retirent après une perte de 20~\% fait
face à une contrainte de \omterm{def:m2:thinking-in-expected-value:frac}{drawdown} que Kelly ignore, et doit dimensionner ses mises à partir du
\omterm{def:m2:thinking-in-expected-value:frac}{risque de ruine}, et non du \omterm{def:m2:thinking-in-expected-value:growth}{taux de croissance}.

\section{Tutoriel~: paris répétés à plusieurs fractions de Kelly}\label{tut:m2:thinking-in-expected-value:sizing}

\begin{tutorial}
\textbf{Objectif.} Calculer la fraction de Kelly et le \omterm{def:m2:thinking-in-expected-value:growth}{taux de croissance}, simuler des paris
répétés à plusieurs multiples de Kelly, vérifier la formule du \omterm{def:m2:thinking-in-expected-value:frac}{risque de ruine}, et trouver le
meilleur multiple lorsque l'edge est estimé.
\textbf{État final~:} \cref{fig:m2:thinking-in-expected-value:growth,fig:m2:thinking-in-expected-value:coin,fig:m2:thinking-in-expected-value:dd,fig:m2:thinking-in-expected-value:uncertain}
et les chiffres du problème du week-end.
\begin{tutsteps}
\item \textbf{Kelly et croissance}, et le simulateur.
\omcode{code/firm/sizing/firm_sizing.py}{14}{41}{La fraction de Kelly, le taux de croissance et les paris répétés.}\label{lst:m2:thinking-in-expected-value:kelly}
\item \textbf{Ruine et erreur d'estimation}~: la probabilité de \omterm{def:m2:thinking-in-expected-value:frac}{drawdown}, son inverse, et le
rétrécissement.
\omcode{code/firm/sizing/firm_sizing.py}{44}{66}{Risque de ruine, fraction pour une tolérance donnée, et rétrécissement pour un edge estimé.}\label{lst:m2:thinking-in-expected-value:ruin}
\item \textbf{Exécuter} \texttt{sizing\_demo.coin\_table()}, \texttt{sizing\_demo.problem()}
et \texttt{fig\_sizing.py}.
\end{tutsteps}
\textbf{Pour aller plus loin.} Rejouer le pile ou face avec le plafond de 250~dollars et comparer des
fractions constantes avec une règle qui mise juste assez pour atteindre le plafond~; puis donner aux
paris des issues à queues épaisses et chercher de combien la meilleure fraction baisse.
\end{tutorial}

\section{Construction~: le module de dimensionnement}\label{bld:m2:thinking-in-expected-value:sizing}

\begin{build}
\bfield{Rôle} Chaque stratégie de la mini-firme est dimensionnée ici~: une fraction du capital à
partir de son edge et de son risque estimés, réduite pour l'erreur d'estimation et pour la tolérance
de la firme aux \omterm{def:m2:thinking-in-expected-value:frac}{drawdowns}.
\bfield{Interface} \texttt{kelly\_fraction(p, b)}~; \texttt{growth(f, p, b)}~;
\texttt{simulate(f, p, n, paths, seed, b, cap)}~; \texttt{drawdown\_probability(c,
alpha)}~; \texttt{fraction\_for\_drawdown(alpha, delta)}~; \texttt{shrinkage(n, sharpe)}~;
\texttt{growth\_uncertain(c, mu, s, n)}.
\bfield{Règles} Richesse multiplicative, mise une fraction fixe de la richesse courante~; le
\omterm{def:m2:thinking-in-expected-value:frac}{risque de ruine} sous sa forme en temps continu~; le rétrécissement dans l'approximation
quadratique~; simulations à graine fixée.
\bfield{Tests d'acceptation} \texttt{code/firm/sizing/tests/}~: Kelly maximise la croissance,
deux fois Kelly a une croissance proche de zéro~; la formule de ruine et son inverse concordent~; les chutes simulées à
moitié correspondent à la formule~; le rétrécissement maximise la croissance espérée.
\bfield{Pour aller plus loin} Kelly pour plusieurs stratégies corrélées (la fraction devient
$\Sigma^{-1}\mu$)~; queues épaisses~; dimensionnement sous une limite de \omterm{def:m2:thinking-in-expected-value:frac}{drawdown} stricte~; mise à jour
bayésienne de l'edge au fil de l'arrivée des transactions (livre~4).
\end{build}

\begin{omsources}
\item V.~Haghani et R.~Dewey, ``Rational decision-making under uncertainty: observed
betting patterns on a biased coin'', octobre 2016.
\item J.~L.~Kelly, ``A new interpretation of information rate'', Bell System Technical
Journal, 1956.
\end{omsources}

\section{Exercices}

\begin{exercise}[$\star$]\label{exo:m2:thinking-in-expected-value:1}
Un pari paie 2 contre 1 et gagne avec une probabilité de 0.4. Quels sont son edge et sa fraction de Kelly~?
\end{exercise}

\begin{exercise}[$\star$]\label{exo:m2:thinking-in-expected-value:2}
Pourquoi la richesse espérée après 300 lancers est-elle un mauvais guide de ce que le joueur finit par avoir~?
\end{exercise}

\begin{exercise}[$\star$]\label{exo:m2:thinking-in-expected-value:3}
Quelle fraction du \omterm{def:m2:thinking-in-expected-value:growth}{taux de croissance} de Kelly complet un demi-Kelly conserve-t-il, dans l'approximation
quadratique~? Et un quart de Kelly~?
\end{exercise}

\begin{exercise}[$\star\star$]\label{exo:m2:thinking-in-expected-value:4}
Avec Kelly complet, quelle est la probabilité de tomber un jour à 20~\% de la richesse initiale~? Avec
un demi-Kelly~?
\end{exercise}

\begin{exercise}[$\star\star$]\label{exo:m2:thinking-in-expected-value:5}
Les investisseurs d'un fonds se retirent s'il perd à un moment 20~\% de son capital initial, et il accepte
une probabilité de 10~\% de subir cette perte. Quel multiple de Kelly doit-il utiliser~?
\end{exercise}

\begin{exercise}[$\star\star$]\label{exo:m2:thinking-in-expected-value:6}
Un edge est estimé avec une $t$-statistique de 2. Quel multiple de la fraction de Kelly estimée
maximise la croissance espérée~?
\end{exercise}

\begin{exercise}[$\star\star\star$]\label{exo:m2:thinking-in-expected-value:7}
\emph{Programmation.} Avec \texttt{simulate} et un plafond de dix fois la mise, estimer la part des
joueurs qui atteignent le plafond en 300 lancers de la pièce à 60~\%, à demi-Kelly et à Kelly,
et comparer avec les 95~\% cités par Haghani et Dewey.
\end{exercise}

\begin{exercise}[$\star\star\star$]\label{exo:m2:thinking-in-expected-value:8}
\emph{Repérer l'erreur.} «~Notre backtest montre un edge avec un ratio de Sharpe de 0.1 par transaction
sur 200 transactions, donc nous misons la fraction de Kelly complète~: elle maximise la croissance.~» Corrigez-le.
\end{exercise}

\section{Problème~: l'edge incertain}

\begin{problem}[{Problème du week-end --- quelle part de Kelly miser}]\label{pb:m2:thinking-in-expected-value:1}
Un trader dispose d'un signal qui, selon lui, gagne 55~\% du temps sur les transactions à \omterm{def:m2:thinking-in-expected-value:edge}{cote} 1 contre 1. Il
a observé 200 transactions et va dimensionner les 500 suivantes comme une fraction du capital.

\medskip\noindent\textbf{Partie I --- L'avantage connu.}
\begin{enumerate}
\item Si le taux de réussite vaut exactement 55~\%, quels sont l'avantage, la fraction de Kelly et le \omterm{def:m2:thinking-in-expected-value:growth}{taux de
croissance} par transaction~?
\item Quels sont la moyenne et l'écart type du profit par unité misée, et le
ratio de Sharpe par transaction~?
\item Avec le Kelly plein, quelle est la probabilité de voir un jour le capital divisé par deux~?
\item Quel multiple de Kelly maintient cette probabilité en dessous de 10~\%~?
\item Pourquoi miser deux fois le Kelly détruit-il l'avantage~?
\end{enumerate}

\medskip\noindent\textbf{Partie II --- L'estimation.}
\begin{enumerate}[resume]
\item Quelle est l'erreur standard de l'avantage sur le taux de réussite estimé à partir de 200 transactions, et sa
statistique $t$~?
\item Quel multiple de la fraction de Kelly estimée maximise la croissance espérée~?
\item Que donne la simulation~?
\item Combien la mise de la fraction de Kelly estimée pleine fait-elle perdre de croissance par rapport au
meilleur multiple~?
\item Au meilleur multiple, quelle est la probabilité de voir un jour le capital divisé par deux~?
\end{enumerate}

\medskip\noindent\textbf{Partie III --- Robustesse.}
\begin{enumerate}[resume]
\item Comment le meilleur multiple évolue-t-il avec un historique de 50 transactions~? Avec 2\,000~?
\item Pourquoi miser trop est-il pire que miser trop peu d'une même quantité~?
\item Que feraient les queues épaisses à la meilleure fraction~?
\item Comment cette réduction se combinerait-elle avec une limite de \omterm{def:m2:thinking-in-expected-value:frac}{drawdown}~?
\item Que dit l'expérience de la pièce sur le dimensionnement en pratique~?
\end{enumerate}

\medskip\noindent\textbf{Partie IV --- Jugement.}
\begin{enumerate}[resume]
\item Pourquoi une société de gestion dimensionnerait-elle ses stratégies au quart de Kelly ou moins~?
\item Comment la fraction doit-elle être mise à jour à mesure que de nouvelles transactions arrivent~?
\item Dans quel cas le cadre de Kelly est-il inadapté~?
\item Énoncer le \emph{résultat nommé}~: le meilleur multiple de Kelly lorsque l'avantage est
estimé à partir de 200 transactions.
\item En une phrase~: pourquoi miser moins que Kelly~?
\end{enumerate}
\end{problem}

\section{Questions d'entretien}

\begin{interviewq}[$\star$ \iqroles{trader, chercheur}]\label{iq:m2:thinking-in-expected-value:1}
Une pièce tombe sur pile dans 60~\% des cas et rapporte la mise à parité. Quelle part de votre
fortune pariez-vous, et pourquoi~?
\end{interviewq}

\begin{interviewq}[$\star$ \iqroles{trader}]\label{iq:m2:thinking-in-expected-value:2}
Quelle est la différence entre maximiser la richesse espérée et maximiser l'espérance
du logarithme de la richesse~?
\end{interviewq}

\begin{interviewq}[$\star\star$ \iqroles{chercheur}]\label{iq:m2:thinking-in-expected-value:3}
Dérivez la fraction de Kelly pour un rendement continu de moyenne $\mu$ et de variance $\sigma^2$.
\end{interviewq}

\begin{interviewq}[$\star\star$ \iqroles{chercheur, risque}]\label{iq:m2:thinking-in-expected-value:4}
Pourquoi les praticiens utilisent-ils le \omterm{def:m2:thinking-in-expected-value:frac}{Kelly fractionnaire}~?
\end{interviewq}

\begin{interviewq}[$\star\star$ \iqroles{trader, chercheur}]\label{iq:m2:thinking-in-expected-value:5}
Vous disposez de deux stratégies non corrélées de ratios de Sharpe 1 et 2. Comment répartissez-vous le
capital entre elles~?
\end{interviewq}

\begin{interviewq}[$\star\star\star$ \iqroles{développeur}]\label{iq:m2:thinking-in-expected-value:6}
Concevez un service de dimensionnement qui fixe chaque jour le capital de chaque stratégie à partir de sa
performance en temps réel, avec des limites.
\end{interviewq}
